\documentclass[book.tex]{subfiles}

\begin{document}
\chapter{Semiconductors}

\label{chap:semiconductors}
\noindent\rule{11cm}{0.4pt} \\
\textbf{Prerequisites:} \autoref{chap:quantum_mechanics}, \autoref{chap:stat_thermo} \\
\textbf{Difficulty Level:} ** \\
\noindent\rule{11cm}{0.4pt}
\vspace{5mm}

Semiconductors are materials with electrical conductivity between that of a metallic conductor and a nonmetallic insulator. Semiconductors have a number of interesting physical properties. For one, their resistivity drops as temperatures rise (unlike metals which have higher resistivity at higher temperatures). Semiconductors can also have their conducting properties altered by the addition of "dopants", impurities introduced into the crystal structure of a semiconductor. A substance with two differently doped regions can form a "semiconductor junction".

In this chapter, we introduce the foundations of semiconductor physics and consider how to build differentiable programs that predict semiconductor structure and behavior. %We also briefly discuss some of the challenges inherent in semiconductor manufacturing.

\section{Basic Semiconductor Physics}

Semiconductors are marked by their unique conductivity behavior, between that for conductors and insulators. To understand this behavior more precisely, we will need to delve into the physics of electronic band structure. An isolated atom has electrons that occupy atomic orbitals. 
%\textcolor{red}{TODO: Is an orbital an eigenfunction of the Schrodinger' equation for the atom?}. 
What happens to the electronic structure then when two atoms come together to form a molecule? Pauli's exclusion principle dictates that two electrons can't have the quantum number in a molecule. This means that each atomic orbital in a diatomic molecule splits into two molecular orbitals %\textcolor{red}{TODO: so does this mean that if each atom had $N$ orbitals, there are $4N$ total or $2N$ at the end?}
% \textcolor{red}{TODO: define quantum number}

Now consider a more complex situation, where $N$ identical atoms come together to form crystal lattice. In this case, the atomic orbitals of individual atoms will overlap. The Pauli exclusion principle then dictates that we can't have two electrons in the solid that have the same quantum numbers. This forces each atomic orbital for each atom to split into $N$ discrete molecular orbitals. Macroscopic solids have very large numbers of atoms, so these $N$ discrete molecular orbitals effectively form a continuum since the energy gap between adjacent molecular orbitals will be very small.

\subsection{Fermi Level}

For a solid state body, the Fermi level is the thermodynamic work required to add one electron to the body. This quantity is usually referred to as $\mu$.

\subsection{Band Gap}

For semiconductors, the band gap is the distance between the valence band and the conduction band. It represents the amount of energy required to promote an electron to participating in conduction.

\begin{figure}
    \centering
    % \includegraphics{figures/Physics/condensed_matter/semiconductors/band_structure.png}
    \includegraphics{figures/Physics/condensed_matter/semiconductors/band gap structure (1).png}
    \caption{The band-gap structure of an example material.}
    \label{fig:bandgap}
\end{figure}
%\textcolor{red}{TODO: Currently sourced from Wiki. Replace with our own figure.}

\subsection{Bravais Lattice and Reciprocal Lattice}

To study the mathematical properties of semiconductor systems, we need to be able to formulate the repeating lattice structure of the material in our equations. Let's start by defining a lattice. Conceptually, a lattice is defined by its unit cell, which is given by three primitive lattice vectors 
\begin{align}
\vec{a_1}, \vec{a_2}, \vec{a_3}
\end{align} 
The Bravais lattice is the infinite set of points described by 

\[ \vec{r} = n_1 \vec{a_1} + n_2 \vec{a_2} + n_3 \vec{a_3}, \quad n_1, n_2, n_3 \in \mathbb{Z} \]

To model physical crystals, we also need to consider the \textit{primitive unit cell}. This unit cell is given by

\[ \textbf{r} = \{x_1 \vec{a_1} + x_2 \vec{a_2} + x_3 \vec{a_3} | 0 \leq x_1, x_2, x_3 < 1 \} \]

We then say that a crystal is an arrangement of atoms within the primitive unit cell. Note that not all the atoms have to be of the same type.

\section{Solving Schrodinger's Equation in Lattices}

Let's consider how to solve Schrodinger's equation for particles in a lattice.  Physically, we expect that the potential energy of the system obeys the following symmetry

\begin{align} 
V(\vec{x} + \vec{a}) = V(\vec{x}) 
\end{align}

That is, the potential energy repeats for each primitive unit cell. The solutions of Schrodinger's equation for such a system are given by Bloch's theorem which states that for electrons in a crystal lattice, with a potential $V$ that obeys the symmetry 

\begin{align} 
V(\vec{x} + \vec{a}) = V(\vec{x}), 
\end{align}

There exists a basis of wave functions which have the following properties: Each wave function $\psi_i$ is an energy eigenstate of Schrodinger's equation. We can write $\psi_i$ as 
\begin{align} 
\psi_i(r) = e^{i k \cdot r} u_i(r)
\end{align}
    
Where $u(r)$ has the same periodicity as the crystal. That is, 
\begin{align}
u(r + \vec{a}) = u(r)
\end{align}
We say that a function with the form
\begin{align}
\psi_{i}(r) = e^{ik\cdot r}u_{i}(r)
\end{align}
is a Bloch wave. Note that the Block theorem means that given bloch state can be identified with a wave vector $k$.
%\subsection{Brillouin Zone}


%\subsection{Free Electron Approximation}

%\textcolor{red}{TODO: Discuss Density Functional Theory for Semiconductors}

%\section{Band Diagrams}

\begin{marginfigure}
   % \includegraphics[width=\linewidth]{figures/Physics/condensed_matter/semiconductors/Pn-junction_zero_bias.png}
   \includegraphics[width=\linewidth]{figures/Physics/condensed_matter/semiconductors/pn junction.png}
   \caption{Band diagram for p–n junction at equilibrium. The depletion region is shaded. }
\label{fig:bandgap}
\end{marginfigure}
% \url{https://commons.wikimedia.org/wiki/File:Pn-junction_zero_bias.png}. TODO: Replace with our own figure


\section{Doping}

Doping is the introduction of impurities into semiconductor for the purpose of changing its properties and behaviors. 

P-type semiconductors have electron acceptors embedded into the underlying semiconductor, whhile N-type semiconductors have electron donors embedded into the underlying semiconductor.


%  \textcolor{red}{TODO: Replace. From \url{https://en.wikipedia.org/wiki/Transistor\#/media/File:MOSFET_Structure.png}}


%\section{Manufacturing Semiconductors}



\end{document}